\subsection{Quem contribui com que parcela numa carteira 3-de-5}
\label{sec:quorum-tres-cinco}
\index{multi-assinatura@multi-assinatura!contagem de parcelas@contagem de parcelas}

Suponhamos que os cinco membros de uma carteira são \(A,B,C,D,E\).
Na manhã em que chegam as contas, só \(A,C,E\) podem assinar. No
caso 2-de-3 era fácil ver à vista desarmada que os dois presentes,
entre si, dispunham de todas as parcelas. Para cinco pessoas vale
a pena seguir a arrumação até ao fim. Agora \(q=5-3+1=3\).
Cada parcela tem três detentores. O número de parcelas diferentes
é \(\binom53=10\), e cada pessoa tem \(\binom42=6\). Nos cinco
registos individuais aparecem, portanto, trinta ocorrências das
mesmas dez parcelas. A chave pública dá a cada parcela o
coeficiente de agregação que lhe corresponde.

Dos três presentes, quem leva para a assinatura cada uma dessas
dez contribuições? Para tornar a contagem transparente, atribuamos
cada parcela ao primeiro dos presentes, por ordem alfabética, que
a tenha. O quadro usa as iniciais dos detentores como índices da
parcela já ponderada \(x_S\).
\begin{center}
\begin{tabular}{c l c}
Signatário & Índices das parcelas que fornece & Total\\ \hline
\(A\) & \(ABC,\ ABD,\ ABE,\ ACD,\ ACE,\ ADE\) & 6\\
\(C\) & \(BCD,\ BCE,\ CDE\) & 3\\
\(E\) & \(BDE\) & 1
\end{tabular}
\end{center}
Cada índice aparece uma vez. O participante \(A\) trouxe as seis
parcelas que conhece. Restam os quatro trios escolhidos só entre
\(B,C,D,E\). Desses quatro, \(C\) conhece três. Sobra \(BDE\),
cujo único detentor presente é \(E\). Se \(E\) tiver de sair antes
de dar a sua contribuição, \(A\) e \(C\) ficam com nove parcelas.
Conhecem o ponto público correspondente à décima, mas a resposta
de assinatura pede a contribuição do respectivo escalar.

Sejam \(s_A,s_C,s_E\) as somas das parcelas nas três linhas.
O quadro verifica a igualdade pública
\[
                   s_AG+s_CG+s_EG=K^s .
\]
Esta escrita permite seguir os pontos sem obrigar pessoa alguma a
reunir os três escalares. Numa resposta de Schnorr de sinal
positivo que use o mesmo desafio \(c\), por exemplo, os termos
\(cs_A,cs_C,cs_E\) contribuem em conjunto com \(cK^s\) depois
da multiplicação por \(G\). As contribuições aleatórias de cada
signatário entram separadamente, como veremos adiante. Se a parcela
\(x_{ACE}\) fosse atribuída a \(E\), retirar-se-ia de \(s_A\) e
somar-se-ia a \(s_E\). A igualdade pública conservar-se-ia. A
atribuição por ordem alfabética dá apenas uma maneira concreta de
fazer a partilha daquela manhã.

Toda a política admite a mesma contagem. Com \(r\) presentes, falta
exactamente uma parcela por cada grupo de \(q\) detentores escolhido
inteiramente entre os \(N-r\) ausentes. O número de parcelas em falta é
\[
              \binom{N-r}{q},
       \qquad \binom{u}{q}=0\ \text{quando }u<q .
\]
Na carteira 3-de-5, com uma pessoa presente faltam \(\binom43=4\).
Com duas falta \(\binom33=1\). Com três, o número cai a zero,
seja qual for o trio. Uma quarta pessoa pode oferecer uma escolha
mais cómoda para as parcelas com vários detentores, e o ponto
\(K^s\) continua a ser o mesmo.

A disponibilidade tem também uma leitura numérica. Tomemos uma
sessão de duração fixada e suponhamos que cada um dos cinco consegue
participar até ao fim com probabilidade \(4/5\), independentemente
dos demais. Para que a sessão fique aquém do limiar, têm de comparecer
zero, uma ou duas pessoas. Com denominador comum \(5^5=3125\), as
três probabilidades têm numeradores
\[
             1,\qquad \binom51\,4=20,\qquad
             \binom52\,4^2=160 .
\]
Assim, a probabilidade de haver um quórum é
\((3125-181)/3125=2944/3125=94{,}208\%\).
O número diz respeito à presença das pessoas sob a hipótese indicada.
Se precisarem todas da mesma ligação e esta estiver disponível com
probabilidade \(9/10\), independente das presenças, a probabilidade
conjunta desce para \(0{,}9\cdot0{,}94208=0{,}847872\). A igualdade
que permite a quaisquer três compor as dez parcelas continua exacta
nos dois cenários. O que mudou foi a hipótese que permite reunir
essas pessoas na mesma sessão.
